% Imporditud: tools/import_eksperimendid.py
% Allikas: Eesti füüsikaolümpiaadi eksperimendi ülesannete
%   kogu (Rael Kalda, 2023), ülesanne Ü151, lahendus L151.
\setAuthor{Erkki Tempel}
\setRound{lõppvoor}
\setYear{2018}
\setNumber{G E1}
\setDifficulty{4}
\setTopic{staatika}
\setVahendid{kaks identset puidust klotsi, koormis massiga m = 100 g, millimeeterpaber, niit, lühike kummipael}
\setPunktid{10}
\setAllikas{Eksperimendi kogu Ü151, lahendus lk 112}
\setLahendusseis{toores}

\prob{Klotsi mass}
Leidke puidust klotsi mass mk.

\hint

\solu
Kinnitame kumminiidi mõlemasse otsa niidid. Ühe niidi otsa kinnitame klotsi külge Hoiame niidi otsast kinni ning venitame klotsiga kumminiidi välja. Laseme kumminiidi lahti ning mõõdame teepikkuse, mille klots läbib. Lisame klotsile 100 grammise raskuse ning venitame uuesti kumminiidi samaplaju välja, kui esimesel juhul ning laseme lahti ja möödame samuti klotsi poolt läbitud teepikkuse. Kumminiidi energia Ep kandub mõlemal juhul üle kotsile ning hõõrdejõud teeb klotsi pidurdamisel tööd A = Fhs = $\mu$mgs. Kuna hõõrdetegur $\mu$ on mõlemal katsel sama, saame avaldada masside suhte esimese ja teise katse jaoks järgmiselt

$m_1$ = $s_2$ $\Rightarrow$ mklots = 100 g $\cdot$ $s_2$ m$^2$ $s_1$ $s_1$ $-$ $s_2$
\probend
